\selectlanguage{ngerman}%
\frenchspacing%
\begin{ledgroupsized}[r]{120mm}%
\footnotesize%
\pstart%
\noindent\textbf{Überlieferung:}%
\pend%
\end{ledgroupsized}%
\begin{ledgroupsized}[r]{114mm}%
\footnotesize%
\pstart%
\parindent -6mm%
\makebox[6mm][l]{\textit{L}}%
Aufzeichnung:
LH XXXVII\;3 Bl.\;62.
Ein Blatt 4\textsuperscript{o},
Ränder beschnitten;
\href{\Scan{LH_37_03_062wz}}{italienisches Papier} (ohne Wasserzeichen).
Zwei ganzseitig beschriebene Seiten;
geringfügige Textverluste durch Beschneidung am linken Rand von Bl.\;\href{\Scan{LH_37_03_062v}}{62\,v\textsuperscript{o}}.
Gestrichene, hier nicht wiedergegebene Entwürfe zu den Diagrammen \bracks{\textit{Fig.\;1}} %
bis \bracks{\textit{Fig.\;3}}.
\pend%
\end{ledgroupsized}%
\vspace{5mm}%
\begin{ledgroup}%
\footnotesize%
\pstart%
\noindent%
\textbf{Datierungsgründe:}
Die Datierung der Aufzeichnung N.\;\ref{60197} beruht auf Eigenschaften des Textträgers sowie auf inhaltlichen Betrachtungen.
Siehe die ausführliche Begründung auf S.\;\pageref{60197_dat}.%
\pend%
\end{ledgroup}%
\selectlanguage{latin}%
\frenchspacing%
\vspace{6mm}%
\normalsize%
\pstart%
\noindent%
\centering%
\bracks{%
\href{\Scan{LH_37_03_062r}}%
{62\,r\textsuperscript{o}}%
}%
\hfill%
\edtext{%
De pondere in plano inclinato%
\indexS{pondus descendens in plano inclinato}%
\indexS{planum inclinatum}
\hfill%
\phantom{\bracks{62\,r\textsuperscript{o}}}%
\linebreak%
per brachium horizonti obliquum%
\indexS{brachium staterae}%
\indexS{statera}%
\indexS{horizon}
in stateram agente%
\indexS{pondus agens in stateram}%
}{%
\lemma{\textit{Am Rand:}}%
\Afootnote{%
Hic casus proponi potest
peritiam mechanicam jactanti.%
\indexS{casus proponendus jactanti peritiam mechanicam}%
\indexS{peritia mechanica}%
}}%
\pend%
\vspace{0.5em}%
\pstart%
\noindent%
\edtext{%
Ponamus
pondera%
\indexS{pondera in statera}
ita esse sita
in
\edtext{%
\textit{H} et \textit{G}%
}{%
\lemma{\textit{H} et \textit{G}}%
\Cfootnote{%
Vgl. die Diagramme \bracks{\textit{Fig.\:1}} bzw. \bracks{\textit{Fig.\:5}}.%
}}
ut%
\bracks{,}%
}{%
\lemma{Ponamus}%
\Bfootnote{%
\textit{(1)}~pondus in %
\textit{(2)}~pondera ita esse sita in\;\textit{H} %
\textit{(a)}~\angles{et}\;\textit{G} %
\textit{(b)}~et\;\textit{G} ut%
~\textit{L}%
}}
\textit{H} descendente ad \textit{S}
in plano inclinato \textit{HS}%
\indexS{planum inclinatum}
et transferente
\edtext{%
staterae%
\indexS{statera}%
}{%
\lemma{staterae}%
\Bfootnote{%
\textit{erg.\;L}%
}}
brachium \textit{AH} in \textit{AS} %
adeoque brachium \textit{AG} in \textit{AI},%
\indexS{brachium staterae}
attollatur pondus perpendiculariter ex \textit{G} in \textit{I}.
Quaeritur
quae debeat esse proportio ponderum%
\indexS{proportio ponderum in statera}
in \textit{H} et \textit{G},
ut sint in aequilibrio.%
\indexS{aequilibrium ponderum in statera}%
\indexS{aequilibrium brachiorum staterae}
\pend%
\pstart%
Ducatur \textit{ST} horizonti parallela,%
\indexS{horizon}
et in eam demittatur perpendicularis \textit{HT}.
Cum
\edtext{%
\bracks{ita}%
}{%
\lemma{ita}%
\Bfootnote{%
\textit{erg. Hrsg.}%
}}
sit,
utique debent esse
pondus in \textit{H},
nempe \astrosun,
ad pondus in \textit{G},%
\indexS{pondus agens in stateram}
nempe $\rightmoon,$
reciproce ut descensus,%
\indexS{descensus brachii staterae}%
\indexS{descensus ponderis in plano inclinato}
id est directe
\edtext{%
\astrosun\ ad $\rightmoon$
ut \textit{GI} ad \textit{HT}.%
}{%
\lemma{\astrosun\ ad $\rightmoon$}%
\Bfootnote{%
\textit{(1)}~ut \textit{GI} ad \textit{HS} %
\textit{(2)}~ut \textit{GI} ad \textit{HT}.%
~\textit{L}%
}}
Porro datur ratio \textit{HT} ad \textit{HS}
quae sit \textit{m}.
Rursus
\edtext{%
ex \textit{S} ducta $S\xi$ normali in \textit{AH},%
}{%
\lemma{ex \textit{S}}%
\Bfootnote{%
\bracks{...} in \textit{AH} %
\textit{erg.\:L}%
}}
datur ratio \textit{HS} ad $\xi S,$
\edtext{%
nempe \textit{n}.
Datur ergo et
ratio
\edtext{%
$ HT :\, \xi S. $%
}{%
\lemma{$ HT :\, \xi S $}%
\Cfootnote{%
Das Verhältnis ist gleich~1.%
}}
Ex \textit{I}%
}{%
\lemma{nempe \textit{n}.}%
\Bfootnote{%
\bars{Datur ergo %
\bracks{...} $ HT :\, \xi S. $ \textit{erg.}} %
\textit{(1)}~Et \textit{I} %
\textit{(2)}~Ex \textit{I}%
~\textit{L}%
}}
agatur $I\beta$
normalis in
\edtext{%
\textit{AG}.
Quia angulus \textit{GAI}
est aequalis angulo \textit{HAS},
erunt triangula $A\xi S$ et $A\beta I$ similia,%
\indexS{triangula similia}
et $ \xi S : \beta I \squaredots AS : AI. $
Id est,
si de initio descensus%
\indexS{descensus brachii staterae}%
\indexS{descensus ponderis in plano inclinato}
agatur,
sunt $ \bracks{AS : AI} \squaredots AH : AG. $
Jam $ \beta I : IG \squaredots AC : AG. $
Ergo
demum componendo
erit%
}{%
\lemma{\textit{AG}.}%
\Bfootnote{%
\textit{(1)}~Erit %
\textit{(a)}~\textit{AS} %
\textit{(b)}~$S\xi : AS$ %
\textit{(c)}~ob %
\textit{(2)}~Ob triangula %
\textit{(3)}~Et %
\textit{(4)}~Quia angulus %
\textit{(a)}~\textit{GIA} %
\textit{(b)}~\textit{GAI} est %
\bracks{...} et $\xi S : \beta I \squaredots AS : AI$
\textit{(aa)}~Jam %
\textit{(bb)}~Id est %
\bracks{...} agatur, sunt %
\textit{(aaa)}~\textit{AI} ad \textit{AS} ut \textit{AG} %
\textit{(bbb)}~$AS : AI$ %
\textit{(ccc)}~\bars{$AS : AI$ \textit{erg. Hrsg.}} $\squaredots\; AH : AG.$ %
\textit{(aaaa)}~Fit ergo %
\textit{(bbbb)}~Jam $\beta I : IG \squaredots AC : AG.$ Ergo demum %
\textit{(aaaaa)}~dantur %
\textit{(bbbbb)}~componendo erit%
~\textit{L}%
}}
$ HT : GI \squaredots mn \cdot \overline{AH \cdot AC} :
\edtext{%
\overline{AG^2} $.
Jam \textit{n},%
}{%
\lemma{$ \overline{AG^2} $.}%
\Bfootnote{%
\bars{Jam \textit{n} \textit{streicht Hrsg.}} Jam \textit{n},%
~\textit{L}%
}}
seu $ HS : \xi S, $
est $ \squaredots\; AH : AD. $
Ergo fiet denique
\edtext{%
\edlabel{LH_37_03_062_062r_nxs-1}%
$ HT : GI \squaredots m \cdot \overline{AH^2 \cdot AC} : \overline{AG^2 \cdot AD}. $
Seu
ut sit aequilibrium%
\indexS{aequilibrium brachiorum staterae}%
\indexS{aequilibrium ponderum in statera}
erit pondus $\rightmoon$ directum
ad pondus \astrosun\ obliquum%
\indexS{pondus agens in stateram}
in ratione composita%
\indexS{ratio composita}
(\textit{m} id est)
sinus anguli inclinationis plani%
\indexS{planum inclinatum}%
\indexS{inclinatio plani}%
\indexS{angulus inclinationis plani}%
\indexS{sinus anguli inclinationis}
ad radium,%
\indexS{radius staterae}
et (\textit{AC} ad \textit{AD}) directa%
\indexS{ratio directa}
perpendicularium%
}{%
\lemma{$ HT : GI $}%
\Bfootnote{%
\textit{(1)}~$\squaredots$ %
\textit{(2)}~seu pondus $\rightmoon$ ad pond %
\textit{(3)}~$\squaredots\ m \cdot \overline{AH^2 \cdot AC} : \overline{AG^2 \cdot AD} $. Seu %
\textit{(a)}~erit p %
\textit{(b)}~ut sit %
\bracks{...} pondus $\rightmoon$ %
\bars{directum \textit{erg.}} ad pondus \astrosun\ %
\bars{obliquum \textit{erg.}} in ratione composita %
\bars{(\textit{m} id est) \textit{erg.}} sinus %
\textit{(aa)}~complementi %
\textit{(bb)}~anguli inclinationis %
\bracks{...} radium, et %
\textit{(aaa)}~\textit{AC} %
\textit{(bbb)}~directa %
\textit{(ccc)}~(\textit{AC} ad \textit{AD}) %
\bars{directa \textit{erg.}} perpendicularium%
~\textit{L}%
}}
ad directiones ex centro,%
\indexS{centrum staterae}
et ($AH^2$ ad $AG^2$) reciproca%
\indexS{ratio reciproca}
duplicata%
\indexS{ratio duplicata}
\edlabel{LH_37_03_062_062r_1a}%
brachiorum.%
\indexS{ratio brachiorum staterae}%
\edlabel{LH_37_03_062_062r_nxs-2}%
\edtext{}{%
\xxref{LH_37_03_062_062r_1a}{LH_37_03_062_062r_1b}%
\lemma{brachiorum.}%
\Bfootnote{%
\bracks{/} %
\textit{(1)}~Idem aliter investigare tentemus %
\textit{(2)}~Verum %
\textit{(3)}~Ducatur jam \textit{IE} parallela ipsi \textit{HS}, erit triangulum \textit{AIE} ad triangulum %
\textit{(4)}~Cum triangula \textit{AIE} et \textit{AHS} sint si %
\textit{(5)}~Ducatur \textit{IE} %
\bracks{...} \textit{AIE} aequalis %
\bars{anguli \textit{ändert Hrsg.}} \textit{AHS}, erunt triangula %
\textit{(a)}~\textit{AHS} %
\textit{(b)}~\textit{AIE} et \textit{AHS} %
\bracks{...} homologorum laterum,%
~\textit{L}%
}}%
\pend%
\begin{edfigure}[2.15em][1.5em]
\hspace{-55mm}%
\includegraphics%
{gesamttex/edit_VIII-4/images/LH_37_03_062_d1a_062r.pdf}%
\edcaption[-0.25em]{%
\hspace{-55mm}%
\edtext{%
\bracks{\textit{Fig.\:1a}}%
}{%
\lemma{\textit{Über dem Diagramm}:}%
\Afootnote[nonum]{%
Figura haec melius in fine paginae versae.%
\textsuperscript{\bracks{a}}%
\indexS{figura}%
\indexS{pagina versa}%
\\[1ex]%
{\footnotesize%
\textsuperscript{\bracks{a}}~%
Figura \bracks{...} versae:
Vgl. das (unvollständige) Diagramm \bracks{\textit{Fig.\:5}}, das in \textit{L} auf Bl.\;\href{\Scan{LH_37_03_062v}}{62\,v\textsuperscript{o}} überliefert ist.%
}%
}%
\lemma{\bracks{\textit{Fig.\:1a}}}%
\Cfootnote[nonum]{%
Das Diagramm entspricht den Angaben im Text.
Die dort nicht erwähnten Konstruktionslinien sind nicht wiedergegeben;
vgl. die Abbildung \bracks{\textit{Fig.\:1b}}.%
}%
}%
}%
\vspace{-10.25em}%
\hspace{65mm}%
\includegraphics[scale=0.20]%
{gesamttex/edit_VIII-4/images/LH_37_03_062_d1b_062r.pdf}%
\edcaption[0.275em]
{%
\hspace{65mm}%
\bracks{\textit{Fig.\:1b}}%
}%
\end{edfigure}
\pstart%
Ducatur \textit{IE} ita
ut sit angulus \textit{AIE} aequalis
angul\bracks{o}
\textit{AHS},
erunt triangula \textit{AIE} et \textit{AHS} similia%
\bracks{,}
at triangula similia%
\indexS{triangula similia}
sunt in ratione duplicata homologorum laterum,%
\indexS{ratio laterum homologorum}%
\indexS{ratio duplicata}%
\indexS{latera triangulorum homologa}%
\edlabel{LH_37_03_062_062r_1b}
ergo triang.\;\textit{AIE},
id est (in initiis) triangulum \textit{AIG},
est ad triangulum \textit{AHS}
ut $AI^2,$
hoc est (in initiis) $AG^2,$
ad $AH^2.$
Sed eadem triangula
\textit{AIG} et \textit{AHS}
sunt ut \textit{AC} in \textit{IG}
ad \textit{AD} in \textit{HS}.
Ergo denique fiet
\textit{IG} ad \textit{HS}
in ratione composita%
\indexS{ratio composita}
$AG^2$ ad $AH^2$ %
\edtext{%
et \textit{AD} ad \textit{AC}.%
}{%
\lemma{et}%
\Bfootnote{%
\textit{(1)}~\textit{A} %
\textit{(2)}~reciproca %
\textit{(3)}~\textit{AD} ad \textit{AC}.%
~\textit{L}%
}}
\edtext{%
Quod consentit priori.%
}{%
\lemma{Quod \bracks{...} priori}%
\Cfootnote{%
Vgl. \refpassage{LH_37_03_062_062r_nxs-1}{LH_37_03_062_062r_nxs-2}.%
}}
\pend%
\pstart%
\edtext{%
\edlabel{LH_37_03_062_062r_add-1}%
Non ergo h\angles{i}c verum est
pondera in plano%
\indexS{planum inclinatum}
inclinat\bracks{o}
hoc modo descendentia%
\indexS{pondus descendens in plano inclinato}
(ut brachia impellant%
\indexS{brachium staterae}%
\indexS{pondus impellens brachium staterae}
quibus affixa non sunt) %
imminuta in ratione%
\indexS{ratio catheti ad hypotenusam}
catheti%
\indexS{cathetum}
ad hypotenusam inclinationis,%
\indexS{hypotenusa inclinationis}
seu potentias,%
\indexS{potentia ponderis in plano inclinato}
esse in reciproca ratione%
\indexS{ratio reciproca}
perpendicularium ex centro \textit{A}%
\indexS{centrum staterae}
ad rectas directionum.
Quod mirum.
Verum esset si
fune ad brachii certum punctum extremum affixo,%
\indexS{funis affixus brachio staterae}%
\indexS{brachium staterae}
potentia traheret%
\indexS{potentia trahens}%
\bracks{;}
ibi enim angulus inclinationis%
\indexS{angulus inclinationis}
inter agendum mutatur.%
\edlabel{LH_37_03_062_062r_add-2}%
}{%
\lemma{Non ergo}%
\Bfootnote{%
\bars{h\angles{i}c \textit{erg.}} verum est %
\textit{(1)}~momenta ponderum in ratione\indexS{momentum ponderis in plano inclinato} %
\textit{(2)}~pondera in plano %
\bars{inclinationis \textit{ändert Hrsg.}} %
\textit{(a)}~hic desc %
\textit{(b)}~hoc modo %
\bracks{...} in ratione %
\textit{(aa)}~inclinationis %
\textit{(bb)}~catheti ad %
\bracks{...} ex centro %
\bars{sphaerae \textit{gestr.}} \textit{A} ad rectas directionum. Quod mirum.
\textit{(aaa)}~Quod tamen %
\textit{(bbb)}~Verum %
\textit{(aaaa)}~esset %
\textit{(bbbb)}~esset %
\textit{(aaaaa)}~cum %
\textit{(bbbbb)}~si fune ad %
\textit{(aaaaa-a)}~immobile %
\textit{(bbbbb-b)}~brachii %
\textit{(aaaaa-aa)}~\angles{cert} %
\textit{(bbbbb-bb)}~certum punctum extremum affixo, potentia %
\textit{(aaaaa-aaa)}~trahit %
\textit{(bbbbb-bbb)}~traheret ibi %
\bracks{...} agendum mutatur. %
\textit{erg.\;L}%
}%
\lemma{Non \bracks{...} mutatur}%
\Cfootnote{%
Die Bemerkung dürfte nach Abfassung des Abschnitts \refpassage{LH_37_03_062_062v_lvd-1}{LH_37_03_062_062v_lvd-2} ergänzt worden sein und bezieht sich vermutlich auf die Ergebnisse von Leibnizens Auseinandersetzung mit %
Fragen der Geostatik aus \textsc{P.\;de Fermat}, \textit{Varia opera mathematica} (Toulouse 1679)\cite{00533} in %
den Exzerpten N.\;\ref{58229} im vorliegenden Band.
Siehe hierüber auch die ausführlichen Datierungsgründe zu N.\;\ref{60197}, S.\;\pageref{60197_dat}.%
}%
}
\bracks{%
\href{\Scan{LH_37_03_062v}}%
{62\,v\textsuperscript{o}}%
}%
\pend%
\pstart%
\edlabel{LH_37_03_062_062v_lvd-1}%
Quoniam
\edtext{%
ostendimus%
}{%
\lemma{ostendimus}%
\Cfootnote{%
Siehe \refpassage{LH_37_03_062_062r_nxs-1}{LH_37_03_062_062r_nxs-2}.%
}}
\edtext{%
esse
$ \rightmoon : \astrosun \squaredots \overline{m \cdot AH^2\cdot AC} : \overline{AG^2 \cdot AD}, $
seu
$ \rightmoon : m\, \astrosun = \overline{AH^2 : AG^2} $
in $ \overline{AC : AD}, $
seu si pondera,%
\indexS{pondera in statera}
in ratione inclinationis%
\indexS{ratio inclinationis}
seu catheti ad hypotenusam%
\indexS{ratio catheti ad hypotenusam}%
\indexS{cathetum}%
\indexS{hypotenusa inclinationis}
imminuta,
sint in ratione perpendicularium%
\indexS{ratio directa}%
\indexS{ratio reciproca}%
\indexS{ratio duplicata}
ex centro staterae%
\indexS{centrum staterae}%
\indexS{statera}
in lineas directionum directa %
et brachiorum duplicata reciproca%
\bracks{,}
erit aequilibrium.%
\indexS{aequilibrium brachiorum staterae}%
\indexS{aequilibrium ponderum in statera}
Jam si coincidant \textit{AH} et \textit{AD},%
}{%
\lemma{esse}%
\Bfootnote{%
\textit{(1)}~$ \astrosun : \rightmoon \squaredots m \cdot AH^2 \cdot A $
\textit{(2)}~$ \rightmoon : \astrosun \squaredots \overline{m\,AH^2\cdot AC} : \overline{AG^2 \cdot AD} $ %
\textit{(a)}~. Jam si coincidant %
\textit{(b)}~, seu $ \rightmoon : m\, \astrosun $ %
\bracks{...} in $ \overline{AC : AD} $ %
\bars{, seu %
\bars{seu \textit{streicht Hrsg.}} %
\bars{si \textit{erg.}} pondera, in %
\bracks{...} imminuta, sint %
\textit{(1)}~hic in aequilibrio %
\textit{(2)}~in %
\textit{(3)}~in ratione %
\bracks{...} erit aequilibrium %
\textit{erg.}%
}. Jam si coincidant \textit{AH} et \textit{AD},%
~\textit{L}%
}}
sitque praeterea planum
non inclinatum%
\indexS{planum inclinatum}
sed perpendiculare,%
\indexS{planum perpendiculare}
seu
\edtext{%
$ m = 1, $
et coincidant denique \textit{AG} et \textit{AC},
fiet%
}{%
\lemma{$ m = 1, $}%
\Bfootnote{%
\textit{(1)}~fiet %
\textit{(2)}~et coincidant %
\bracks{...} \textit{AC}, fiet%
~\textit{L}%
}%
\lemma{coincidant \bracks{...} \textit{AG} et \textit{AC}}%
\Cfootnote{%
Vgl. das Diagramm \bracks{\textit{Fig.\;2}.}%
}%
}
$ \rightmoon : \astrosun \squaredots AD : AC, $
quod verum esse scimus
\edtext{%
aliunde.%
}{%
\lemma{aliunde}%
\Cfootnote{%
Gemeint ist das Hebelgesetzt.
Siehe \textsc{Archimedes},
\textit{De planorum aequilibriis}~%
I, prop.\;6--7.\cite{01010}%
}}
\edlabel{LH_37_03_062_062r_rmo-1}%
Quodsi
\edtext{%
faciamus \textit{AH} parallelam horizonti,%
\indexS{horizon}
erit
$ HS : HT \squaredots AH : AD, $
hoc est $ m = AD : AH .$
Quod si ergo praeterea 
$ AG = AC, $
his valoribus substitutis%
\indexS{valor substitutus}
\ledpb%
fiet:%
}{%
\lemma{faciamus}%
\Bfootnote{%
\textit{(1)}~$ AG = AC, $ et %
\textit{(2)}~\textit{AH} parallelam %
\bracks{...} $ HS : HT \squaredots AH : AD $ %
\textit{(a)}~. Jam si praeterea %
\textit{(b)}~, hoc est %
\bracks{...} praeterea  $ AG = AC, $ %
\textit{(aa)}~fiet
\textit{(bb)}~his valoribus substitutis fiet:%
~\textit{L}%
}}
$ \rightmoon : \astrosun = \overline{AD \cdot AH \cdot AC} : AC^2 \cdot AD, $
id est
$ \rightmoon : \astrosun \squaredots AH : AC. $%
\edlabel{LH_37_03_062_062r_rmo-2}
\edtext{%
Quod
\edlabel{LH_37_03_062_062v_dfs-1}%
paradoxum.%
\indexS{paradoxum}%
}{%
\lemma{Quod paradoxum}%
\Cfootnote{%
Siehe die Erläuterung zu \refpassage{LH_37_03_062_062r_add-1}{LH_37_03_062_062r_add-2}.%
}}%
\edtext{}{%
\xxref{LH_37_03_062_062v_dfs-1}{LH_37_03_062_062v_dfs-2}%
\lemma{paradoxum.}%
\Bfootnote{%
\textit{(1)}~Ideo amplius examinandum. Triangulum \textit{AGI} est ad triangulum \textit{AHS} (in initiis) ut $AG^2$ ad $AH^2,$ et eadem sunt ut \textit{AC} in \textit{GI} ad %
\textit{(2)}~\bracks{/} Ideo amplius exami %
\textit{(3)}~%
\bars{Ideo amplius examinandum. \textit{erg.}} \bracks{/} %
Video esse verum,
\textit{(a)}~nam %
\textit{(b)}~nam in casu aequilibrii %
\textit{(aa)}~est \astrosun\ ad $\rightmoon,$ ut %
\textit{(aaa)}~\textit{HT} a %
\textit{(bbb)}~\angles{in} %
\textit{(ccc)}~\textit{C\angles{I}} %
\textit{(bb)}~sunt pondera %
\bars{ad \textit{erg. Hrsg.}} initia descensuum %
\bars{vel ascensuum \textit{erg.}} reciproce, seu %
\bracks{...} ad \textit{HT}.%
~\textit{L}%
}}
Ideo amplius examinandum.%
\edlabel{LH_37_03_062_062v_lvd-2}%
\pend%
\pstart%
Video esse verum,
nam in casu aequilibrii%
\indexS{casus aequilibrii}%
\indexS{aequilibrium brachiorum staterae}%
\indexS{aequilibrium ponderum in statera}
sunt pondera
\edtext{%
\bracks{ad}%
}{%
\lemma{ad}%
\Bfootnote{%
\textit{erg. Hrsg.}%
}}
initia descensuum%
\indexS{descensus brachii staterae}%
\indexS{descensus ponderis in plano inclinato}
vel ascensuum%
\indexS{ascensus brachii staterae}%
\indexS{ascensus ponderis in statera}
reciproce,
seu \astrosun\ ad $\rightmoon$
est
ut \textit{CI} ad \textit{HT}.%
\edlabel{LH_37_03_062_062v_dfs-2} %
Jam secet \textit{AS} ipsam \textit{HT} in $\omega,$
patet
triangula $AH\omega$ et \textit{AGI} esse similia.%
\indexS{triangula similia}
Ergo erit
$ CI : H\omega \squaredots AC : AH. $
\pend%
\pstart%
Verum in casu initii
coincidunt $H\omega$ et \textit{HT},
nam
\edtext{%
differentia%
}{%
\lemma{differentia}%
\Bfootnote{%
\textit{erg.\;L}%
}}
$\omega T$ est incomparabiliter minor%
\indexS{differentia incomparabiliter minor}
ipsa $H\omega.$
Habemus ergo
esse
$ CI : HT \squaredots AC : AH. $
Ergo denique \astrosun\ ad $\rightmoon$
ut \textit{AC} ad \textit{AH},
\edtext{%
ut ante.%
}{%
\lemma{ut ante}%
\Cfootnote{%
Siehe \refpassage{LH_37_03_062_062r_rmo-1}{LH_37_03_062_062r_rmo-2}.%
}}
\pend%
\pstart%
Et id
ipsum redditur manifestius,
si pondus \astrosun%
\indexS{pondus descendens in plano inclinato}
ponatur sub %
plano inclinato,%
\indexS{planum inclinatum}
patet enim
utroque modo rem eodem
\edtext{%
redi\bracks{re},%
}{%
\lemma{redit,}%
\Bfootnote{%
\textit{L\;ändert Hrsg.}%
}}
et revera
in isto descensu%
\indexS{descensus ponderis in plano inclinato}%
\indexS{descensus brachii staterae}
pondus \astrosun%
\indexS{pondus descendens in plano inclinato}
non sustineri
a plano inclinato,%
\indexS{planum inclinatum}
etsi ita videatur,
sed retineri tantum
ne inferius decurrat
in ipsum brachium \textit{AS} continuatum,%
\indexS{brachium staterae}
ultra quam oportet,
seu ut maneat in recta illa inclinata.
\pend%
\begin{edfigure}[2.0em][0.00em]%
\hspace{6mm}%
\includegraphics%
{gesamttex/edit_VIII-4/images/LH_37_03_062_d2_062v.pdf}%
\edcaption[-0.5em]{%
\hspace{5.5mm}%
\bracks{\textit{Fig.\:2}}}
\end{edfigure}
\begin{edfigure}[2.25em][0.00em]%
\hspace{-72mm}%
\includegraphics%
{gesamttex/edit_VIII-4/images/LH_37_03_062_d3_062v.pdf}
\edcaption[0.5em]{%
\hspace{-72mm}%
\bracks{\textit{Fig.\:3}}%
}%
\vspace{-9.00em}%
\hspace{74mm}%
\includegraphics%
{gesamttex/edit_VIII-4/images/LH_37_03_062_d4_062v.pdf}
\edcaption[0.5em]{%
\hspace{74mm}%
\bracks{\textit{Fig.\:4}}%
}%
\end{edfigure}
\newpage%
\begin{edfigure}[0.00em][1.00em]
\includegraphics%
{gesamttex/edit_VIII-4/images/LH_37_03_062_d5_062v.pdf}
\edcaption[-0.5em]{%
\lbrack\textit{Fig.\:5}\rbrack}%
\end{edfigure}
\vspace{5.0em}%
\pstart%
\noindent%
\centering%
*     *     *
\pend%
\vspace{5.0em}%
\selectlanguage{german}%
\frenchspacing%
\footnotesize%
\pstart%
\noindent%
\textls{Ausführliche Datierungsgründe:}%
\pagelabel{60197_dat}
Die titellose und undatierte Aufzeichnung N.\;\ref{60197} ist auf einem Träger überliefert, der zwar keine Spuren eines Wasserzeichens aufweist, anhand seiner äußeren Merkmale aber als römisches\indexO{Rom} Papier identifiziert werden kann.
In der Farbe und der Anordnung der Stege ähnelt er nämlich den Trägern weiterer im vorliegenden Band edierter Texte (etwa N.\;\ref{58656-A}, N.\;\ref{58258}, N.\;\ref{58218}, N.\;\ref{58240}, N.\;\ref{58220} und N.\;\ref{58235}), deren römische\indexO{Rom} Herkunft feststellbar ist.
Auf gleichem Papier sind ebenfalls Leibnizens Auszüge mit Bemerkungen aus \textsc{P.\;de Fermat}, \textit{Varia opera mathematica} (Toulouse 1679)\cite{00533} überliefert, welche zum Teil unter N.\;\ref{58229} im vorliegenden Band ediert sind.
Eine Ergänzung in N.\;\ref{60197} (\refpassage{LH_37_03_062_062r_add-1}{LH_37_03_062_062r_add-2}) deutet vermutlich auf diese Auszüge hin, insbesondere auf Leibnizens Auseinandersetzung mit Fragen der Geostatik, die im Mittelpunkt des Briefwechsels zwischen Pierre de Fermat\indexN{\textls{Fermat}, Pierre de 1601\textendash1665} und Gilles Personne de Roberval\indexN{\textls{Roberval} (Robervallius), Gilles Personne de 1602\textendash1675} standen; siehe N.\;\ref{58229} im vorliegenden Band.
Die Aufzeichnung N.\;\ref{60197} dürfte somit nicht vor den Auszügen mit Bemerkungen N.\;\ref{58229} 
entstanden sein, jedenfalls -- wie diese letzteren -- zu der Zeit, als Leibniz sich in Rom\indexO{Rom} aufhielt.
Auch N.\;\ref{60197} ist demgemäß auf den Zeitraum von Mitte April bis November 1689 zu datieren.
Eine frühere Entstehung ist ausgeschlossen, eine spätere unwahrscheinlich.
\pend%
\selectlanguage{latin}%
\frenchspacing%
\normalsize%
